\section{问题定义与设计目标}

\subsection{BLoP 基础模式与升级目标}

本文将 BLoP 作为基础可信审计模式。该模式通过可信计算、可信 Agent、异常事件监测和 PBFT 风格共识形成完整信任链，其默认假设是由可信节点集合对关键审计事件进行充分确认。对于安全等级较高、审计频率较低或责任要求极强的场景，全节点确认模式具有直接价值。本文关注的是另一个工程问题：当审计事件频率升高、链上记录数量增多、不同事件风险等级差异明显时，是否可以在不改变 BLoP 可信治理主线的前提下，增加一条轻量化执行路径，使中低风险事件通过可信委员会和批量存证完成确认，高风险事件仍可回退到全节点模式。

\subsection{问题定义}

考虑一个由多个机构参与的数据流通场景。数据提供方将数据对象切分为若干分片，并将分片或计算证据保存在链下存储节点中；审计节点周期性检查分片完整性或证据一致性；联盟链用于记录数据承诺、审计摘要、异常处理和节点信誉变化。系统希望在不暴露原始数据、不把大体积证据直接上链的前提下，完成分布式数据校验。

本文关注的问题可以表述为：在链下数据分片或审计证据可能损坏、丢失或被恶意修改的条件下，如何设计一种低链上开销的分布式校验机制，使系统能够完成异常发现、数据恢复、责任记录和信誉更新，并在联盟链中形成可追溯证据。

\subsection{威胁模型}

本文假设联盟链参与方经过准入管理，但其中部分节点可能出现故障或拜占庭行为。链下存储节点可能删除分片、返回错误分片、延迟响应或提交不一致证据；审计节点可能不按期执行审计、伪造审计状态或与存储节点合谋。联盟链共识层容忍不超过 $f$ 个拜占庭节点，并要求委员会中可投票节点数满足 $2f+1$ 的确认阈值。本文暂不处理密码学协议本身被攻破、所有高信誉节点同时合谋、可信硬件侧信道攻击和业务语义造假等强攻击。

\subsection{设计目标}

TrustFlowAudit 的设计目标包括：

\begin{itemize}
  \item \textbf{低链上开销：}链上只保存哈希承诺、Merkle 根、审计状态、恢复状态和信誉变化，不保存原始数据和大体积证明。
  \item \textbf{过程可追溯：}审计请求、异常发现、恢复处理和节点惩罚均形成链上摘要记录，便于后续仲裁和复核。
  \item \textbf{恢复闭环：}当损坏分片不超过纠删码恢复阈值时，系统应能触发恢复并更新新的数据承诺。
  \item \textbf{委员会轻量确认：}基于可信评分选择小规模审计委员会，为 BLoP 全节点确认模式提供轻量化可选路径。
  \item \textbf{可回退全节点模式：}当数据对象风险较高、委员会可信度不足或审计争议较大时，系统可以回退到 BLoP 式全可信节点确认。
  \item \textbf{边界清晰：}本文方法不替代 BLoP 的完整可信计算链路，也不替代完整 PDP/PoR 密码学协议；这些能力可作为基础能力或后续扩展模块接入。
\end{itemize}
